\section{Common domination and the information in exact positions}
\label{sec:v26-domination}

Three properties must be separated: domination of a statistical family by
one sigma-finite measure, absolute continuity with respect to a designated
reference member, and mutual absolute continuity of the members.  Moving
support need not prevent the first property.  The distinction also separates
the scalar-coordinate local experiments from the richer position protocol.

\begin{proposition}[Common reference measures]
\label{prop:v26-common-domination}
The family \eqref{eq:v22-thinned-family} is dominated, for every \(n\), by
\(\mu=\operatorname{Leb}|_U+\delta_\dagger\).  Its censored family is
dominated by \(\operatorname{Leb}|_U+\delta_\dagger+\delta_\partial\).
The finite products have the corresponding product dominating measures.
The original members need not be absolutely continuous with respect to
\(P_{n,0}\); the common-collar construction establishes this stronger
property for the censored members and their censored reference.

On the fixed finite strip \(B=D\times[-R,R]\), let
\[
 d\Lambda_z=\rho(u,v)\mathbf1_{\{y>U(u,v)z\}}\,du\,dv\,dy,
 \qquad d\Lambda_* =\rho(u,v)\,du\,dv\,dy.
\]
Then the Poisson law \(\mathsf P_* =\operatorname{PPP}(\Lambda_*)\) is a
common dominating probability measure for the entire finite-strip family,
and
\begin{equation}\label{eq:v26-poisson-envelope}
 \frac{d\mathsf P_z}{d\mathsf P_*}(N)
 =\exp\{\Lambda_*(B)-\Lambda_z(B)\}
           \prod_{q\in N}\mathbf1_{\{y(q)>U(q)z\}}.
\end{equation}
The empty product is one.  Multiplication by a fixed batch proportion gives
the same assertion for each batch of \eqref{eq:v23-et-poisson-intensity}.
\end{proposition}

\begin{proof}
The first statements follow by writing the displayed Lebesgue densities
and atomic masses.  They do not involve division by the reference density.
For example, on \([0,2]\),
\[
 f_\theta(x)=\frac{2(1+\theta-x)_+}{(1+\theta)^2},\qquad 0\le\theta<1,
\]
is a family of probability densities, but
\(P_\theta((1,1+\theta))=\theta^2/(1+\theta)^2>0\) for \(\theta>0\),
whereas \(P_0((1,1+\theta))=0\).

For the Poisson assertion write \(A_z=\{y>Uz\}\cap B\).  Under
\(\mathsf P_*\), restrictions to \(A_z\) and its complement are
independent Poisson processes.  The probability of no point in the
complement is \(\exp\{-\Lambda_*(B\setminus A_z)\}\).  Multiplying its
indicator by the reciprocal of that probability has expectation one and
leaves the independent process on \(A_z\) unchanged.  The resulting law
has intensity \(\Lambda_*|_{A_z}=\Lambda_z\), proving
\eqref{eq:v26-poisson-envelope}.  The argument includes the zero-intensity
and empty-configuration cases.
\end{proof}

A zero-shift Poisson member need not dominate the others.  If \(Uz<0\) on
a set of positive weighted endpoint measure, there is positive intensity
in \(\{Uz<y<0\}\), where the reference has none.  The event of at least
one point there has positive alternative probability and zero reference
probability.  This is a failure of reference absolute continuity, not of
common domination.  It also excludes a Gaussian shift with mutually
absolutely continuous members as the local limit in such a direction.
The common-collar Gaussian limit and the moving-ceiling Poisson limit remain
different experiments for the reasons supplied by their boundary rates and
support events, not by an incorrect classification of their dominating
measures.

For a fixed compact parameter set, choose \(R\) beyond all displacements.
The positive tail outside the strip is parameter independent and the negative
tail is empty.  Adjoining the common positive-tail process preserves the
common domination just proved, also for the projective experiment of
Remark~\ref{rem:v24-projective-poisson}.  This makes no uniform assertion
about an unbounded parameter set without a common displacement bound.

\subsection{Exact geometric positions and scalar coarsening}

There are genuine support-singular position experiments.  The following
comparison is local and registered; it is not a claim about every compact
global class in the reconstruction theorem.

\begin{proposition}[A strict position--scalar comparison]
\label{prop:v26-position-deficiency}
Consider a persistent anchored analytic contact family in which the contact
point and tangent are fixed, but one obstacle varies through distinct
homothetic copies about that point.  Keep a nondegenerate facing channel and
a fixed admissible even-flight design.  Conditional on success, suppose the
first endpoint has a density with respect to boundary arclength, as in the
positive flux model.  Then distinct members' raw position laws are mutually
singular.  On any uncountable interval of such copies, the raw position
family has no common sigma-finite dominating measure.

The scalar endpoint laws obtained by the fixed transverse projection are,
in contrast, continuous in total variation as the homothety parameter
varies.  If \(\mathsf E^{(k)}_{\rm sc}\) and
\(\mathsf E^{(k)}_{\rm pos}\) denote \(k\) successful scalar and position
records on a fixed parameter interval, then
\begin{equation}\label{eq:v26-position-deficiency}
 \delta(\mathsf E^{(k)}_{\rm sc},\mathsf E^{(k)}_{\rm pos})\ge\tfrac12,
 \qquad k<\infty.
\end{equation}
This conditional-success statement does not discard the preparation cost
in a theorem about an unconditional acquisition transcript.
\end{proposition}

\begin{proof}
Let \(K\) be the convex obstacle with anchor \(p\in\partial K\).  Along
each ray from \(p\) entering its interior, strict convexity gives exactly
one other boundary point.  Scaling about \(p\) multiplies its distance
from \(p\).  Thus the boundaries of two distinct positive homothetic
copies meet only at \(p\).  A continuous arclength-density endpoint hits
\(p\) with probability zero.  The first-position marginal, and hence the
complete position record, is supported on disjoint Borel curve sets after
\(p\) is removed.  This proves mutual singularity.

A sigma-finite measure has at most countably many disjoint measurable sets
of positive measure: intersect with its countable finite-measure cover and
then group intersections with measure at least \(1/m\).  A common dominating
measure would assign positive measure to each of the uncountably many
full-probability curve sets, a contradiction.

The tangent and anchor are fixed along this family, so transverse projection
is one parameter-independent map.  In the persistent small contact window,
the finite generating action and the flux depend smoothly on the homothety
parameter.  After integrating residual time, the scalar density has the
form \(A_\theta(d-E_{J,\theta})_+\) on one fixed bounded box with a positive
normalizer.  Uniform continuity of the factors and
\(|x_+-y_+|\le|x-y|\) give total-variation continuity.  The same proof
applies to the limiting factorized law.

For any two parameters and any Markov kernel \(K\), the triangle inequality
and contraction of total variation imply
\[
 \|P^{\rm pos}_0-P^{\rm pos}_1\|_{\rm TV}
 \le 2\sup_\theta\|KP^{\rm sc}_\theta-P^{\rm pos}_\theta\|_{\rm TV}
                 +\|P^{\rm sc}_0-P^{\rm sc}_1\|_{\rm TV}.
\]
This applies to the \(k\)-fold laws.  Their position distance is one; their
scalar distance tends to zero when two distinct parameters approach each
other, since product total variation is at most \(k\) times the one-record
distance.  Infimizing over kernels gives
\eqref{eq:v26-position-deficiency}.
\end{proof}

The example does not require circular symmetry.  A strictly convex analytic
oval with support function
\[
 h(\theta)=1+\tfrac1{50}\cos(2\theta)+\tfrac1{100}\sin(3\theta)
\]
has \(h+h''\ge43/50>0\).  Its nonzero harmonic orders two and three
exclude a nonidentity rotational symmetry.  Fix a boundary anchor and use
nearby homothetic copies, with a facing obstacle kept at positive gap.
A sufficiently small parameter interval preserves the local channel.
This constructs the stated local comparison without asserting the global
cycle and spanning-tree hypotheses for an unspecified periodic table.

Exact injectivity of the transverse \emph{law} map in
Theorem~\ref{thm:v26-single-offset-global} is compatible with this finite-sample
comparison.  Distribution-level determination and statistical equivalence
of experiments are different assertions.  The former is the intrinsic inverse
proved here; the latter is asserted only for the explicitly matched local
experiments in the deficiency theorems.
